\documentclass[12pt]{article}
\usepackage{geometry}
\geometry{top=15mm, bottom=20mm, left=15mm, right=15mm}
\usepackage{booktabs}
\usepackage{array}
\usepackage{tabularx}
\usepackage{hyperref}
\usepackage{parskip}
\usepackage{setspace}
\usepackage{titlesec}
\usepackage{enumitem}
\usepackage{fancyhdr}
\usepackage{tikz}
\usepackage{amsmath, amssymb}
\usepackage{graphicx}
\usepackage{float}
\usepackage{listings}
\usepackage{xcolor}
\usepackage{caption}
\usepackage{subcaption}
\usepackage{longtable}
\usepackage{verbatim}
\usetikzlibrary{positioning, calc}

\newcolumntype{Y}{>{\raggedright\arraybackslash}X}

\setlength{\headheight}{15pt}
\pagestyle{fancy}
\fancyhf{}
\rhead{CAP5610 -- Machine Learning}
\lhead{Project Final Report Team 4}
\rfoot{\thepage}

\titleformat{\section}{\large\bfseries}{\thesection.\;}{0em}{}[\titlerule]
\titleformat{\subsection}{\normalsize\bfseries}{\thesubsection\;}{0em}{}

\definecolor{codebg}{RGB}{245,245,245}
\definecolor{codecomment}{RGB}{106,153,85}
\definecolor{codekey}{RGB}{0,0,200}
\definecolor{codestr}{RGB}{163,21,21}
\lstset{
backgroundcolor=\color{codebg},
basicstyle=\ttfamily\footnotesize,
breaklines=true,
breakatwhitespace=false,
columns=flexible,
keepspaces=true,
commentstyle=\color{codecomment},
keywordstyle=\color{codekey}\bfseries,
stringstyle=\color{codestr},
showstringspaces=false,
frame=single,
rulecolor=\color{black!30},
numbers=left,
numberstyle=\tiny\color{black!50},
xleftmargin=2em,
framexleftmargin=1.5em,
captionpos=b
}

\hypersetup{
colorlinks=true,
linkcolor=black,
urlcolor=blue,
citecolor=blue
}

\title{Project Final Report Team 4\\[0.5em]\large CAP 5610 Machine Learning}
\author{
    Project Title: Predicting Yelp Review Star Ratings: A Comparative Study of Traditional and Deep Learning Methods\\
    Team Name/Number: \underline{\hspace{3in}}
}
\date{May 4, 2026}

\begin{document}
\sloppy

\begin{center}
    {\LARGE \textbf{Machine Learning Term Project Team 4}} \\[0.5em]
    {\Large Final Report} \\[1em]
    {\large Predicting Yelp Review Star Ratings: \\
    A Comparative Study of Traditional and Deep Learning Methods} \\[1em]
    Alexander Nardi, Anthony Mahon, Benedetto Falin, Keating Sane, Mohammed Mamdouh \\[0.4em]
    May 4, 2026
\end{center}

\setcounter{tocdepth}{2}
\tableofcontents
\newpage

\section{Team Information}

\subsection{Team Members and Roles}

\subsection{Responsibility Distribution}

\subsection{Team Assignment Table}
\begin{table}[H]
    \centering
    \begin{tabularx}{\textwidth}{p{0.36\textwidth} Y}
        \toprule
        \textbf{Team Member}                  & \textbf{Models}                     \\
        \midrule
        Alexander Nardi (Project Coordinator) & Linear SVM, LSTM                    \\
        Anthony Mahon                         & Kernel SVM, Transformer encoder     \\
        Benedetto Falin                       & Naive Bayes, Multi Layer Perceptron \\
        Keating Sane                          & Decision Tree/Random Forest, CNN    \\
        Mohammed Mamdouh                      & Logistic Regression, LLM            \\
        \bottomrule
    \end{tabularx}
\end{table}


\newpage
\section{Problem Setup}

\subsection{Classification Task}


\subsection{Dataset}


\subsection{Dataset Size and Splits}

\begin{table}[h]
    \centering
    \begin{tabular}{lll}
        \toprule
        Split      & Samples & Notes \\
        \midrule
        Training   &         &       \\
        Validation &         &       \\
        Testing    &         &       \\
        \bottomrule
    \end{tabular}
\end{table}

\subsection{Input and Output Format}


\subsection{Class Labels}

\begin{table}[h]
    \centering
    \begin{tabular}{lll}
        \toprule
        Label ID & Class Name & Meaning \\
        \midrule
        0        &            &         \\
        1        &            &         \\
        2        &            &         \\
        3        &            &         \\
        4        &            &         \\
        \bottomrule
    \end{tabular}
\end{table}

\subsection{Basic Dataset Statistics}


\subsection{Project Goals}


\section{Implementation Description}

\subsection{Repository and Codebase Overview}

\subsubsection{Repository Structure}

The codebase is organized such that each model is in its own directory and utilities that are shared across several models are found in \texttt{utils/}:

\vspace{0.5em}

\begin{lstlisting}[language=bash]
CAP5610TermProject/
|-- Linear_SVM/
|-- Kernel_SVM/
|-- Naive_Bayes/
|-- Decision_Tree/
|-- Random_Forest/
|-- Logistic_Regression/
|-- LSTM/
|-- MultiLayer_Perceptron/
|-- CNN/
|-- Transformers/
|-- Large_Language_Model/
|-- utils/
|-- requirements.txt
`-- README.md
\end{lstlisting}

\subsubsection{Shared Utilities}

The \texttt{utils/} package exports a common API that all models scripts can utilize:

\vspace{0.5em}

\begin{lstlisting}[language=bash]
utils/
|-- args.py                    # common CLI flags
|-- data.py                    # HuggingFace yelp_review_full loader + split
|-- embeddings.py              # GloVe / fastText loaders for CNN
|-- evaluation.py              # metrics, plots, results logging
|-- randomness.py              # DEFAULT_SEED, set_random
|-- text_features.py           # shared TF-IDF vectorizer
|-- timer.py                   # timed_step context manager
`-- tuning.py                  # Optuna driver + best-config persistence
\end{lstlisting}

\subsubsection{Model Specific Directories}

Each model's directory contains its script, a tuning log, a results log, a saved best-config, and a confusion matrix for the best final result. For example, \texttt{CNN/}:

\vspace{0.5em}

\begin{lstlisting}[language=bash]
CNN/
|-- CNN.py                     # runnable model script
|-- confusion_matrix_cnn.png   # confusion matrix from best final run
|-- best_config.json           # winning hyperparameters from tuning
|-- results_log.md             # every runs metrics, params, runtime
`-- tuning_log.md              # every Optuna trials params + score
\end{lstlisting}

\subsection{Data Preparation and Experimental Pipeline}

\subsubsection{Data Loading Workflow}

All models read the dataset through a single entry point: \texttt{utils.data.load\_yelp\_data()}, which calls \texttt{datasets.load\_dataset("yelp\_review\_full")}. The dataset downloads on first use and is cached locally. Following runs then read from cache.

\subsubsection{Training/Validation/Test Protocol}

The official split of the dataset is 650,000 training reviews and 50,000 test reviews. The loader then takes an additional validation split (by default 10\%) out of the training partition using \texttt{sklearn.model\_selection.train\_test\_split}. The test partition is never touched during tuning. Validation runs use a default 150,000-review subsample of training data for fast iteration; final runs use the full 650,000-review training set (unless specified otherwise for some models) and evaluate on the held-out 50,000-review test partition (triggered by \texttt{--final}).

\subsubsection{Text Preprocessing}

Sparse-feature models (Logistic Regression, Naive Bayes, Decision Tree, Random Forest, Linear/Kernel SVM, MLP) consume TF--IDF features built directly from raw review text by \texttt{TfidfVectorizer}, which lowercases, tokenizes, and applies \texttt{strip\_accents="unicode"} internally. Sequence models (CNN, LSTM, Transformer) lowercase the text and replace punctuation with whitespace before tokenizing on whitespace, then map tokens to fixed-size index sequences (CNN: \texttt{max\_len = 256};  LSTM: \texttt{max\_len = 512}; Transformer: \texttt{max\_len = 160}). Tokens outside of the vocabulary collapse to \texttt{<unk>} and zero-padded positions are masked at the embedding layer. The decoder-only LLM (default \texttt{distilgpt2}) bypasses this pipeline entirely and instead applies the model's pretrained HuggingFace \texttt{AutoTokenizer} (subword BPE) directly to the raw text, padding/truncating to \texttt{max\_length = 256}.

\subsubsection{Feature Engineering}

\textbf{TF--IDF:} \texttt{fit\_tfidf\_features()} fits a single \texttt{TfidfVectorizer} with \texttt{sublinear\_tf=True} (dampens the contribution of repeated words within a review via $1 + \log(\text{tf})$) and four tunable parameters: \texttt{max\_features} (cap on the vocabulary size), \texttt{ngram\_range} (whether to include bigrams in addition to unigrams), \texttt{min\_df} (drop terms appearing in fewer than this many documents, removing rare/typo tokens), and \texttt{max\_df} (drop terms appearing in more than this fraction of documents, removing extremely common words). Final feature dimensionalities range from 10,000 (MLP) to 40,000 (Decision Tree, Random Forest, Logistic Regression).

\textbf{Word vocabulary:} Each word-level sequence model builds its own top $N$ vocabulary from the training text with index 0 reserved for the padding token (\texttt{<pad>}) and index 1 reserved for the unknown word token (\texttt{<unk>}). CNN: 25,000 words (fixed); LSTM: 30,000 words (tuned over 10k--75k); Transformer: 15,000 words (tuned over 15k--30k). Reviews are truncated or zero-padded to the model's \texttt{max\_len}.

\textbf{Embeddings:} The CNN supports either a randomly initialized 256-dimensional embedding lookup (trained from scratch) or one of several pretrained sources: GloVe 6B 100d/300d, GloVe 42B 300d, GloVe 2024 WikiGigaword 300d, and fastText wiki-news subword 300d, which can all be selected with mutually exclusive CLI flags. Pretrained matrices are loaded into the embedding layer at model construction time. The LSTM uses a trainable \texttt{nn.Embedding} lookup tuned over {128, 256, 300} dimensions with no pretrained option. The Transformer learns both token and positional embeddings from scratch (token embedding tuned over {128, 256} dimensions, positional embedding sized to \texttt{max\_len}). The decoder-only LLM gets its subword embeddings from the pretrained \texttt{distilgpt2} checkpoint and tunes them jointly with the rest of the model.

\subsubsection{Randomness and Reproducibility}

A single global seed (\texttt{DEFAULT\_SEED = 0}) is set at the start of every run via \texttt{set\_random\_seed()}, which seeds all of the following: Python's \texttt{random}, NumPy, and Torch (including CUDA), and pins \texttt{torch.backends.cudnn.deterministic = True} with \texttt{benchmark = False}. The same seed is passed into the dataset splitter and into every Optuna study, so re-running any script reproduces the same train/validation split and the same trial order.

\subsubsection{Hyperparameter Tuning Workflow}

Tuning is invoked with the \texttt{--tune} flag and uses Optuna Bayesian optimization through \texttt{utils.tuning.tune\_model()}. Each model defines an \texttt{objective()} that samples a configuration, builds features, trains on a subset of the training data ranging from roughly 6,000 reviews (LLM) to 650,000 reviews (CNN, Linear SVM) depending on the model's compute cost per trial, and returns macro-F1 on the validation split. Trial budgets range from 6 (Logistic Regression) to 30 (CNN), with most models in the 15--20 range. When tuning completes, \texttt{save\_best\_config()} writes the winning configuration plus metadata (seed, embedding source, train size, validation split, trial count, best macro-F1) to \texttt{best\_config.json}, and the tuning helper appends a Markdown summary of all trials to \texttt{tuning\_log.md}. Following runs load that JSON automatically unless \texttt{--default} is passed. Naive Bayes is the exception: \texttt{--tune} runs a manual alpha sweep over a fixed grid rather than an Optuna study, and writes \texttt{nb\_alpha\_sweep.png} instead of a tuning log.

\subsection{Traditional Machine Learning Models}

\subsubsection{Logistic Regression}
\textbf{Folder/Script:} \texttt{Logistic\_Regression/Logistic\_Regression.py}\\
\textbf{Team Member:}

\paragraph{Model Description}


\paragraph{Input Representation}


\paragraph{Training Procedure}


\paragraph{Hyperparameters}


\paragraph{Tuning Strategy}


\paragraph{Implementation Notes}


\subsubsection{Naive Bayes}
\textbf{Folder/Script:} \texttt{Naive\_Bayes/naive\_bayes.py}\\
\textbf{Team Member:}

\paragraph{Model Description}


\paragraph{Input Representation}


\paragraph{Training Procedure}


\paragraph{Hyperparameters}


\paragraph{Tuning Strategy}


\paragraph{Implementation Notes}


\subsubsection{Decision Tree}
\textbf{Folder/Script:} \texttt{Decision\_Tree/Decision\_Tree.py}\\
\textbf{Team Member:} Keating Sane

\textbf{Model Description:} A single decision tree classifier from \texttt{sklearn.tree.DecisionTreeClassifier}. The tree is built top-down by repeatedly choosing one TF--IDF feature and a threshold that best separates the five star-rating classes (using Gini impurity), then recursing on each side. Predictions follow the chain of feature-vs-threshold tests from the root to a leaf, where the majority class is returned. Implemented as a baseline against the bagged ensemble (Random Forest).

\textbf{Input Representation:} TF--IDF vectors produced by the shared \texttt{fit\_tfidf\_features()} helper with \texttt{sublinear\_tf=True}, \texttt{max\_features = 40,000}, \texttt{ngram\_range = (1, 2)}, \texttt{min\_df = 10}, \texttt{max\_df = 0.9}. No additional preprocessing.

\textbf{Training Procedure:} Final runs fit a single tree on the full 650,000-review training partition with the tuned config and evaluate on the 50,000-review test partition. Validation runs use a 150,000-review training subsample with a 10\% validation split for fast iteration. CPU-only.

\textbf{Hyperparameters:} The tunable parameters are \texttt{max\_depth} (maximum tree depth, which caps model capacity and the rate at which the tree can carve up the feature space), \texttt{min\_samples\_leaf} (minimum training rows per leaf, which prevents the tree from creating tiny leaves that memorize individual reviews), and the four TF--IDF parameters listed above. The splitting criterion is sklearn's default Gini impurity.

\textbf{Tuning Strategy:} Optuna Bayesian optimization with 20 trials on a 200,000-review tuning subset (10\% validation split). Search space: \texttt{tfidf\_features} $\in$ \{5k, 10k, 20k, 40k\}, \texttt{ngram\_max} $\in$ \{1, 2\}, \texttt{min\_df} $\in$ \{3, 5, 10\}, \texttt{max\_df} $\in$ \{0.85, 0.9, 0.95\}, \texttt{max\_depth} $\in$ \{50, 100, 150, 300, None\}, \texttt{min\_samples\_leaf} $\in [1, 10]$. Macro-F1 on the validation split was the objective.

\textbf{Implementation Notes:} Tree fitting on the 200k tuning subset takes well under a minute per trial on CPU, but the full 650k-review final fit grows to roughly an hour as the tree depth scales. With the best tuned settings, the model memorizes most of the training set: train accuracy is far above the macro-F1 it achieves on validation/test. The model was kept as an interpretable baseline to compare with the ensemble.

\subsubsection{Random Forest}
\textbf{Folder/Script:} \texttt{Random\_Forest/Random\_Forest.py}\\
\textbf{Team Member:} Keating Sane

\textbf{Model Description:} Ensemble of decision trees from \texttt{sklearn.ensemble.RandomForestClassifier}. Each tree is fit on a bootstrap sample of the training rows with a random subset of features considered at every split (\texttt{max\_features = "sqrt"}); class predictions are averaged across the forest. The randomized feature subsampling reduces the variance that the single Decision Tree exhibits.

\textbf{Input Representation:} TF--IDF vectors from \texttt{fit\_tfidf\_features()} with \texttt{sublinear\_tf=True}, \texttt{max\_features = 40,000}, \texttt{ngram\_range = (1, 2)}, \texttt{min\_df = 10}, \texttt{max\_df = 0.85}. Identical pipeline to the single tree, with a slightly tighter \texttt{max\_df}.

\textbf{Training Procedure:} Final runs fit a 500-tree forest on the full 650,000-review training partition with the tuned config and evaluate on the 50,000-review test partition. Validation runs use a 150,000-review training subsample with a 10\% validation split. Tree fitting is parallelized across cores via \texttt{n\_jobs = -1}; CPU-only.

\textbf{Hyperparameters:} The tunable parameters are mirrored from the Decision Tree, with two additions: \texttt{n\_estimators} (number of trees in the forest) and \texttt{max\_features} (number of features randomly considered at each split).

\textbf{Tuning Strategy:} Optuna Bayesian optimization with 20 trials on a 200,000-review tuning subset (10\% validation split). Search space mirrors the Decision Tree's parameters and adds \texttt{n\_estimators} $\in$ \{100, 200, 300, 500\} and \texttt{max\_features} $\in$ \{"sqrt", "log2"\}. Macro-F1 on the validation split was the objective.              

\textbf{Implementation Notes:} A 500-tree forest on the 200k tuning subset fits in roughly three minutes per trial with \texttt{n\_jobs = -1}; the full 650k-review final fit grows to about 44 minutes. Bootstrap sampling plus the \texttt{max\_features = "sqrt"} per-split feature subset force each tree to see a different slice of the data and consider a different subset of words at every split, so the trees make different mistakes and the variance that the single Decision Tree exhibits averages out.

\subsubsection{Linear SVM}
\textbf{Folder/Script:} \texttt{Linear\_SVM/Linear\_SVM.py}\\
\textbf{Team Member:}

\paragraph{Model Description}


\paragraph{Input Representation}


\paragraph{Training Procedure}


\paragraph{Hyperparameters}


\paragraph{Tuning Strategy}


\paragraph{Implementation Notes}


\subsubsection{Kernel SVM}
\textbf{Folder/Script:} \texttt{Kernel\_SVM/Kernel\_SVM.py}\\
\textbf{Team Member:}

\paragraph{Model Description}


\paragraph{Input Representation}


\paragraph{Training Procedure}


\paragraph{Hyperparameters}


\paragraph{Tuning Strategy}


\paragraph{Implementation Notes}


\subsection{Deep Learning Models}

\subsubsection{Multi-Layer Perceptron}
\textbf{Folder/Script:} \texttt{MultiLayer\_Perceptron/ml\_perceptron.py}\\
\textbf{Team Member:}

\paragraph{Model Description}


\paragraph{Input Representation}


\paragraph{Architecture}


\paragraph{Training Procedure}


\paragraph{Hyperparameters}


\paragraph{Tuning Strategy}


\paragraph{Implementation Notes}


\subsubsection{Convolutional Neural Network}
\textbf{Folder/Script:} \texttt{CNN/CNN.py}\\
\textbf{Team Member:} Keating Sane

\textbf{Model Description:} A Kim (2014) style TextCNN implemented in PyTorch. An embedding layer maps token indices to dense vectors, several 1D convolutional filters of different kernel widths run in parallel over the embedded sequence (acting as learned $n$-gram detectors), max-over-time pooling collapses each feature map to a single value per filter, the pooled outputs are concatenated, passed through dropout, and projected to five class logits via a fully connected layer. Trained end-to-end with cross-entropy loss.

\textbf{Input Representation:} Reviews are lowercased, punctuation is replaced with whitespace, and the text is tokenized on whitespace. The training corpus produces a 25,000-word vocabulary with index 0 reserved for the padding token (\texttt{<pad>}) and index 1 reserved for the unknown-word token (\texttt{<unk>}). Each review is mapped to a 256-token index sequence (truncated or zero-padded). The embedding source is chosen at the CLI: either a randomly initialized 256-dimensional lookup (default) or one of five pretrained options: GloVe 6B 100d, GloVe 6B 300d, GloVe 42B 300d, GloVe 2024 WikiGigaword 300d, or fastText wiki-news subword 300d. Pretrained matrices are loaded directly into the embedding layer at model construction.

\textbf{Architecture:} Embedding ($V \times E$) $\to$ parallel \texttt{Conv1d} filters at multiple kernel widths $\to$ ReLU $\to$ max-over-time pooling per filter $\to$ concatenate pooled outputs $\to$ dropout $\to$ fully connected layer producing 5-class logits. Padding positions are masked at the embedding via \texttt{padding\_idx = 0}. Optimizer is Adam; gradients are clipped to a max L2 norm of 1.0.

\textbf{Training Procedure:} Final runs train on the full 650,000-review training partition and evaluate on the 50,000-review test partition. Validation runs use a 150,000-review training subsample with a 10\% validation split. Each epoch ends with a held-out validation pass that tracks the lowest-loss checkpoint; training stops early if validation loss has not improved for 3 consecutive epochs, and the lowest-loss checkpoint is restored before evaluation. Auto-selects MPS, CUDA, or CPU at startup.

\textbf{Hyperparameters:} The tunable parameters are \texttt{num\_filters} (number of feature maps per kernel width), \texttt{filter\_sizes} (the set of kernel widths run in parallel, wider filters capture longer $n$-grams), \texttt{dropout} (regularization on the concatenated feature vector before the fully connected layer), \texttt{lr} (Adam learning rate), \texttt{batch\_size}, and \texttt{epochs} (effective via early stopping). The embedding source is selected at the CLI rather than tuned; the choice is recorded in \texttt{best\_config.json}'s metadata so following runs can detect a mismatch.

\textbf{Tuning Strategy:} Optuna Bayesian optimization with 30 trials on the full 650,000-review training partition (10\% validation split). Search space: \texttt{num\_filters} $\in$ \{100, 200, 300\}, \texttt{filter\_sizes} $\in$ \{"3,4,5", "2,3,4,5"\}, \texttt{dropout} $\in [0.2, 0.6]$, \texttt{lr} $\in [10^{-4}, 10^{-2}]$ (log scale), \texttt{batch\_size} $\in$ \{64, 128\}, \texttt{epochs} $\in [5, 15]$. Macro-F1 on the validation split was the objective. The best epoch identified by early stopping is written back into \texttt{best\_config.json} so final runs reproduce that training duration.

\textbf{Implementation Notes:} Training is GPU-bound; on an RTX 5090 a final fit on 650k reviews completes in under a minute per epoch. Because the embedding source materially affects training, the saved choice is recorded in \texttt{best\_config.json}'s metadata so following runs use the same embeddings the config was tuned with.

\subsubsection{Transformer (Encoder Only)}
\textbf{Folder/Script:} \texttt{Transformers/Transformers.py}\\
\textbf{Team Member:}

\paragraph{Model Description}


\paragraph{Input Representation}


\paragraph{Architecture}


\paragraph{Training Procedure}


\paragraph{Hyperparameters}


\paragraph{Tuning Strategy}


\paragraph{Implementation Notes}


\subsubsection{Large Language Model (Decoder Only)}
\textbf{Folder/Script:} \texttt{Large\_Language\_Model/Large\_Language\_Model.py}\\
\textbf{Team Member:}

\paragraph{Model Description}


\paragraph{Pretrained Backbone}


\paragraph{Tokenization and Input Format}


\paragraph{Training Procedure}


\paragraph{Hyperparameters}


\paragraph{Tuning Strategy}


\paragraph{Implementation Notes}


\subsection{Cross-Model Engineering Choices}

\subsubsection{Common CLI Design}

Every model script imports \texttt{common\_parser()} from \texttt{utils.args}, which registers four shared flags: \texttt{--tune} (run hyperparameter search), \texttt{--final} (full training set, evaluate on the test partition), \texttt{--discard} (run without writing any artifacts), and \texttt{--default} (ignore any saved \texttt{best\_config.json} and use hard-coded defaults). Model-specific flags (e.g.\ the CNN's embedding-source selectors, the SVMs' \texttt{--size}, the LLM's checkpoint and resume flags) are layered on top of this shared parser. \texttt{--tune} is mutually exclusive with \texttt{--final} and \texttt{--default}.

\subsubsection{Logging and Artifact Saving}

Run output is structured by helpers in \texttt{utils.evaluation}: \texttt{print\_run\_header()} prints the model name and contextual info; \texttt{print\_value\_section()} prints the parameter table; \texttt{timed\_step()} prints elapsed time for each step during loading/training. After evaluation, \texttt{save\_results()} appends a Markdown section to \texttt{results\_log.md} with accuracy, macro-precision, macro-recall, macro-F1, runtime, parameters, and metadata; the same helper guards against overwriting a better run by comparing macro-F1 against the existing entry. For final runs, \texttt{plot\_confusion\_matrix()} writes a PNG figure into the model folder. \texttt{--discard} bypasses all writes.

\subsubsection{Hardware/Runtime Considerations}

Sparse feature models run on CPU and parallelize through \texttt{n\_jobs=-1} where supported. The CNN, LSTM, Transformer, MLP, and LLM auto-select MPS, CUDA, or CPU at startup and use a fixed size token batch loader. Final run hardware varied: NVIDIA RTX 5090 (Linear SVM, CNN), AMD Ryzen 7 9800X3D (Decision Tree, Random Forest), AMD Ryzen 7 7800X3D (Naive Bayes, MLP), Tesla T4 (Transformer), and Intel i7-1165G7 (Kernel SVM). Final runtimes are recorded per-run in each \texttt{results\_log.md}.

\section{Usage Guide}

This section explains how a user can setup our project to reproduce the results in this report.

\subsection{Environment Setup}

\subsubsection{Prerequisites}

Python 3.11 or newer, \texttt{pip}, and roughly 5 GB of free disk space for the cached dataset and pretrained embedding files. A CUDA or MPS capable GPU is recommended for the CNN, Transformer, and LLM but is not required. All other models run on CPU.

\subsubsection{Installation}

\begin{lstlisting}[language=bash]
# clone
git clone <repo-url> && cd CAP5610TermProject

# create env (Python 3.11)
python -m venv .venv && source .venv/bin/activate

# install deps
pip install -r requirements.txt
\end{lstlisting}

\subsubsection{Dataset Access}

The dataset is fetched on first run via \texttt{datasets.load\_dataset("yelp\_review\_full")} from HuggingFace. No credentials or manual downloads are required; following runs read from the local HuggingFace cache. Pretrained GloVe and fastText files used by the CNN download on demand the first time a corresponding embedding flag is passed.

\subsection{Common Run Modes}

\begin{longtable}{p{0.18\textwidth} p{0.32\textwidth} p{0.38\textwidth}}
    \toprule
    Flag               & Purpose & Notes \\
    \midrule
    \texttt{--tune}    & Run the Optuna hyperparameter search and save the best config. & Trains on a model-specific subset and evaluates on the validation split. \\
    \texttt{--final}   & Train on the full 650,000-review training partition (unless otherwise specified for some models) and evaluate on the 50,000 review test partition. & Loads the saved \texttt{best\_config.json} unless \texttt{--default} is passed. Writes a confusion matrix PNG. \\
    \texttt{--discard} & Skip writing any artifacts (results log, tuning log, best config, confusion matrix). & Useful for debugging. \\
    \texttt{--default} & Ignore saved \texttt{best\_config.json} and use the hard-coded defaults in the script. & Useful when rerunning with a different embedding source than the one the saved config was tuned with (e.g.\ CNN), or when comparing tuned vs.\ untuned baselines. \\
    \bottomrule
\end{longtable}


\subsection{How to Run Each Model}

Every model is invoked with \texttt{python <Directory>/<Script>.py} plus the desired flag combination. A typical workflow is to run \texttt{--tune} once to populate \texttt{best\_config.json}, then run \texttt{--final} to produce the test-set numbers reported in Section~\ref{sec:results}.

\textbf{Inputs:} raw review text loaded automatically from the HuggingFace cache; for the CNN, an optional pretrained embedding flag.\\[0.5em]
\textbf{Outputs:} a printed run header, parameter table, timings, evaluation metrics, and (for final runs) a confusion matrix.\\[0.5em]
\textbf{Artifacts produced:} an appended section in the models \texttt{results\_log.md}; for tuning runs, an appended section in \texttt{tuning\_log.md} and a refreshed \texttt{best\_config.json}; for final runs, \texttt{confusion\_matrix\_*.png}.

\vspace{0.5em}

\begin{lstlisting}[language=bash]
# Hyperparameter search (one time per model)
python Logistic_Regression/Logistic_Regression.py --tune
python Naive_Bayes/naive_bayes.py --tune
python Decision_Tree/Decision_Tree.py --tune
python Random_Forest/Random_Forest.py --tune
python Linear_SVM/Linear_SVM.py --tune
python Kernel_SVM/Kernel_SVM.py --tune
python MultiLayer_Perceptron/ml_perceptron.py --tune
python LSTM/LSTM.py --tune
python CNN/CNN.py --tune --glove-2024-wikigiga
python Transformers/Transformers.py --tune
python Large_Language_Model/Large_Language_Model.py --tune

# Final test-set evaluation using the saved best config
python Logistic_Regression/Logistic_Regression.py --final
python Naive_Bayes/naive_bayes.py --final
python Decision_Tree/Decision_Tree.py --final
python Random_Forest/Random_Forest.py --final
python Linear_SVM/Linear_SVM.py --final
python Kernel_SVM/Kernel_SVM.py --final
python MultiLayer_Perceptron/ml_perceptron.py --final
python LSTM/LSTM.py --final
python CNN/CNN.py --final --glove-2024-wikigiga
python Transformers/Transformers.py --final
python Large_Language_Model/Large_Language_Model.py --final

# Debugging: validation-mode run without writing anything
python CNN/CNN.py --discard --default
\end{lstlisting}

\textbf{Additional flags.} A few models layer extra options on top of the four shared flags:
\begin{itemize}
    \item \textbf{CNN} (mutually exclusive embedding source selectors): \texttt{--glove-6b-100d}, \texttt{--glove-6b-300d}, \texttt{--glove-42b}, \texttt{--glove-2024-wikigiga}, \texttt{--fasttext-wiki-subword}. Omitting all of them uses a randomly initialized 256-dimensional embedding lookup.
    \item \textbf{Linear SVM} and \textbf{Kernel SVM}: \texttt{--size N} overrides the training subset size (defaults: 650{,}000 for Linear SVM, 10{,}000 for Kernel SVM).
    \item \textbf{Large Language Model}: \texttt{--model-name NAME} (default \texttt{distilgpt2}) chooses the decoder-only base model; \texttt{--smoke} runs a tiny model on a tiny subset for sanity checks; \texttt{--train-size N} overrides the non-final training subset size; \texttt{--trials N} sets the number of Optuna trials in \texttt{--tune} mode (default 8); \texttt{--resume} resumes from \texttt{Large\_Language\_Model/saved\_model} and \texttt{--resume-from PATH} resumes from an arbitrary checkpoint. \texttt{--tune} additionally blocks \texttt{--resume}/\texttt{--resume-from}, and \texttt{--final} blocks \texttt{--smoke}.
\end{itemize}

\subsection{Expected Output Files}

\begin{longtable}{p{0.18\textwidth} p{0.36\textwidth} p{0.34\textwidth}}
    \toprule
    Type                     & Location & Purpose \\
    \midrule
    Best configuration       & \texttt{<Model>/best\_config.json} & Hyperparameters and metadata from the most recent tuning run. \\
    Tuning logs              & \texttt{<Model>/tuning\_log.md} & Markdown record of every Optuna trial, noting parameters, val macro-F1, and the runtime. Naive Bayes is the exception: it runs a manual alpha sweep producing \texttt{nb\_alpha\_sweep.png} instead. \\
    Results logs             & \texttt{<Model>/results\_log.md} & Markdown record of every best validation and final run: accuracy, macro-precision/recall/F1, parameters, metadata, runtime. \\
    Confusion matrix  & \texttt{<Model>/confusion\_matrix\_*.png} & PNG figure written on final runs. \\
    \bottomrule
\end{longtable}

\section{Evaluation Metrics}

\subsection{Notation}


\subsection{Accuracy}
\textbf{Formula:}
\[
    \text{}
\]
\textbf{Interpretation:}

\subsection{Macro Precision}
\textbf{Formula:}
\[
    \text{}
\]
\textbf{Interpretation:}

\subsection{Macro Recall}
\textbf{Formula:}
\[
    \text{}
\]
\textbf{Interpretation:}

\subsection{Macro F1-Score}
\textbf{Formula:}
\[
    \text{}
\]
\textbf{Interpretation:}

\subsection{Confusion Matrix}
\textbf{Definition/Formula:}
\[
    \text{}
\]
\textbf{Interpretation:}

\subsection{Why These Metrics Fit This Project}


\section{Results}
\label{sec:results}

\subsection{Validation Experiments}


\subsection{Final Experiments}


\subsection{Hardware and Software Environment}

\begin{table}[h]
    \centering
    \begin{tabular}{ll}
        \toprule
        Item          & Value \\
        \midrule
        OS            &       \\
        Python        &       \\
        CPU           &       \\
        GPU           &       \\
        Key Libraries &       \\
        \bottomrule
    \end{tabular}
\end{table}

\subsection{Tuning Budget}


\subsection{Experiment Tracking}


\subsection{Comparison Table}

\begin{longtable}{p{0.18\textwidth} p{0.08\textwidth} p{0.1\textwidth} p{0.1\textwidth} p{0.1\textwidth} p{0.1\textwidth} p{0.1\textwidth} p{0.14\textwidth}}
    \toprule
    Model                  & Split & Accuracy & Macro Precision & Macro Recall & Macro F1 & Runtime & Notes \\
    \midrule
    Logistic Regression    &       &          &                 &              &          &         &       \\
    Naive Bayes            &       &          &                 &              &          &         &       \\
    Decision Tree          &       &          &                 &              &          &         &       \\
    Random Forest          &       &          &                 &              &          &         &       \\
    Linear SVM             &       &          &                 &              &          &         &       \\
    Kernel SVM             &       &          &                 &              &          &         &       \\
    Multi-Layer Perceptron &       &          &                 &              &          &         &       \\
    CNN                    &       &          &                 &              &          &         &       \\
    Transformer            &       &          &                 &              &          &         &       \\
    Large Language Model   &       &          &                 &              &          &         &       \\
    \bottomrule
\end{longtable}

\subsection{Per-Model Result Summaries}

\subsubsection{Logistic Regression}


\subsubsection{Naive Bayes}


\subsubsection{Decision Tree}


\subsubsection{Random Forest}


\subsubsection{Linear SVM}


\subsubsection{Kernel SVM}


\subsubsection{Multi-Layer Perceptron}


\subsubsection{Convolutional Neural Network}


\subsubsection{Transformer}


\subsubsection{Large Language Model}


\subsection{Visualizations}

\subsubsection{Confusion Matrices}


\subsubsection{Tuning Curves/Loss Curves/Accuracy Curves}


\subsection{Comparative Discussion}

\subsubsection{Best-Performing Models}


\subsubsection{Worst-Performing Models}


\subsubsection{Traditional ML vs. Deep Learning}


\subsubsection{Performance vs. Compute Tradeoffs}


\subsubsection{Error Patterns}


\section{Conclusions \& Future Work}

\subsection{Key Findings}


\subsection{Lessons Learned}


\subsection{Limitations}


\subsection{Future Improvements}


\subsection{Final Takeaway}

\end{document}
